\documentclass[a4paper,12pt]{article}
\usepackage[T2A]{fontenc}
\usepackage[utf8]{inputenc}
\usepackage[russian]{babel}
\usepackage{amsmath,amssymb,mathtools}
\renewcommand{\familydefault}{\sfdefault}
\usepackage{icomma}
\usepackage[a4paper,margin=22mm,top=23mm,bottom=22mm]{geometry}
\usepackage{microtype}
\usepackage{xcolor}
\usepackage{enumitem}
\usepackage{tikz}
\usetikzlibrary{arrows.meta,positioning,shapes.geometric,calc}
\usepackage{tabularx,booktabs}
\usepackage{fancyhdr}
\usepackage{setspace}
\usepackage{hyperref}
\definecolor{accent}{HTML}{008080}
\definecolor{darkaccent}{HTML}{004D4D}
\definecolor{textgray}{HTML}{333333}
\hypersetup{colorlinks=true,linkcolor=darkaccent,pdftitle={Чек-лист. Экстремумы в прикладных задачах},pdfauthor={Лёвин Артём Александрович}}
\setlength{\parindent}{0pt}
\setlength{\parskip}{5pt}
\setstretch{1.12}
\setlist[itemize]{leftmargin=1.55em,itemsep=3pt,topsep=4pt}
\setlist[enumerate]{leftmargin=1.7em,itemsep=4pt,topsep=4pt}
\newcommand{\sectionrule}{\par\vspace{2pt}{\color{accent}\rule{\linewidth}{0.6pt}}\vspace{4pt}}
\newcommand{\lessonsection}[1]{\section*{\color{accent}#1}\sectionrule}
\newcommand{\lessonstep}[1]{\par\vspace{5pt}{\color{accent}\bfseries #1}\par\vspace{-3pt}}
\newcommand{\result}[1]{\begingroup\setlength{\fboxsep}{4pt}\fcolorbox{darkaccent}{white}{$\displaystyle #1$}\endgroup}
\pagestyle{fancy}
\fancyhf{}
\fancyhead[L]{\footnotesize\color{textgray}Экстремумы в прикладных задачах}
\fancyhead[R]{\footnotesize\color{textgray}Тимофей Васильченко}
\fancyfoot[C]{\footnotesize\thepage}
\renewcommand{\headrulewidth}{0.55pt}
\renewcommand{\headrule}{\hbox to\headwidth{\color{accent}\leaders\hrule height \headrulewidth\hfill}}
\setlength{\headheight}{24pt}
\begin{document}
\thispagestyle{empty}
\begin{center}
\vspace*{18mm}
{\color{accent}\small\bfseries МАТЕМАТИКА}
\vspace{20mm}

{\Large Чек-лист}
\vspace{6mm}

{\color{accent}\LARGE\bfseries Экстремумы в прикладных задачах}
\vspace{2mm}

{\color{accent}\large Производная и исследование функции}
\vspace{16mm}

Подготовка к ЕГЭ по профильной математике
\vspace{25mm}

{\Large Тимофей Васильченко}
\vspace{16mm}

13 ноября 2026 г.
\end{center}
\vfill
\begin{center}
Лёвин Артём Александрович

эксклюзивно для Тимофея Васильченко
\vspace{8mm}

\begin{tikzpicture}[x=1cm,y=1cm]
  \draw[accent,line width=.8pt] (0,0) .. controls (1,0.48) and (1.75,-0.48) .. (2.8,0)
  .. controls (3.65,0.4) and (4.5,-0.3) .. (5.6,0);
\end{tikzpicture}
\end{center}
\clearpage

\lessonsection{1. Главный принцип: что требуется исследовать}
В прикладной задаче слова «наибольшее» и «наименьшее» указывают на поиск экстремума некоторой величины. Сначала выразите эту величину через одну переменную, а затем исследуйте функцию.

\lessonstep{Что установить до вычислений}
\begin{itemize}
\item \textbf{Целевая функция:} величина, которую нужно сделать максимальной или минимальной (мощность $P$, плотность распределения $f$, расстояние $l$).
\item \textbf{Переменная:} что реально может изменяться ($R$, $v$, $\alpha$).
\item \textbf{Параметры:} фиксированные числа и физические постоянные ($E$, $r$, $M$, $T$, $h$, $u$).
\item \textbf{Допустимые значения:} ограничения, заданные математикой и физическим смыслом.
\end{itemize}

\lessonstep{Полезные формулы дифференцирования}
Для дифференцируемых $U=U(x)$, $V=V(x)$ при $V\ne 0$:
\[
 (UV)'=U'V+UV',\qquad \left(\frac{U}{V}\right)'=\frac{U'V-UV'}{V^2}.
\]
Сложная функция и элементарные производные:
\[
 (e^{g(x)})'=g'(x)e^{g(x)},\quad ((g(x))^2)'=2g(x)g'(x),
\]
\[
 (\sin x)'=\cos x,\qquad (\cos x)'=-\sin x.
\]

\lessonstep{От уравнения $F'(x)=0$ к ответу}
\begin{enumerate}
\item Упростите производную так, чтобы можно было установить её знак.
\item Найдите критические точки и исключите недопустимые значения.
\item Определите знак $F'$ на интервалах. $F'>0$ означает возрастание, $F'<0$ --- убывание.
\item Докажите, что найденная точка действительно даёт требуемый экстремум, и вычислите искомую величину.
\end{enumerate}

\textbf{Важно.} Условие $F'(x)=0$ само по себе не доказывает существование максимума или минимума. При необходимости исследуйте границы допустимого интервала.
\clearpage

\begin{center}
{\color{accent}\Large Алгоритм: поиск экстремума в прикладной задаче}
\end{center}
\vspace{8mm}
\begin{center}
\begin{tikzpicture}[
  >=Latex, node distance=10mm,
  block/.style={draw=accent!75,rounded corners=2pt,align=center,text width=11.1cm,minimum height=12mm,inner sep=4pt,font=\small},
  decision/.style={diamond,draw=accent!75,aspect=3.8,align=center,text width=6.4cm,inner sep=2pt,font=\small},
  arr/.style={->,line width=.7pt,draw=accent}
]
\node[block] (goal) {Выделите величину, для которой\\требуется максимум или минимум};
\node[block,below=of goal] (formula) {Запишите её как $F(x)$; отделите переменную от констант};
\node[block,below=of formula] (domain) {Укажите допустимую область $D$\\и физические ограничения};
\node[block,below=of domain] (diff) {Найдите $F'(x)$, упростите дроби\\и вынесите общие множители};
\node[decision,below=13mm of diff] (choice) {Есть допустимые решения $F'(x)=0$?};
\node[block,below=17mm of choice] (signs) {Расставьте знаки $F'$ на области $D$, исследуйте границы};
\node[block,below=of signs] (answer) {Докажите нужный экстремум, вычислите\\величину и укажите единицы};
\foreach \p/\q in {goal/formula,formula/domain,domain/diff,diff/choice,signs/answer} \draw[arr] (\p)--(\q);
\draw[arr] (choice.south) -- node[right,font=\small] {да} (signs.north);
\draw[arr] (choice.west) -| node[pos=.21,above,font=\small] {нет} ($(signs.west)+(-.75,0)$) -- (signs.west);
\end{tikzpicture}
\end{center}
\clearpage

\lessonsection{2. Максимальная мощность на сопротивлении}
\textbf{Условие.} Источник с ЭДС $E=60$ В и внутренним сопротивлением $r=0{,}8$ Ом питает резистор сопротивлением $R$. Найдите $R>0$, при котором выделяемая на резисторе мощность наибольшая.

\lessonstep{Шаг 1. Вывести целевую функцию}
Сила тока и тепловая мощность:
\[
 I=\frac{E}{r+R},\qquad P=I^2R.
\]
Подставляем $I$ в формулу мощности:
\[
 P(R)=\frac{E^2R}{(r+R)^2},\qquad R>0.
\]
Здесь переменной является только $R$; величины $E$ и $r$ фиксированы.

\lessonstep{Шаг 2. Найти и упростить производную}
Используем правило производной частного:
\begin{align*}
 P'(R)&=\frac{E^2(r+R)^2-E^2R\cdot2(r+R)}{(r+R)^4}\\[2pt]
 &=\frac{E^2(r+R)\bigl((r+R)-2R\bigr)}{(r+R)^4}\\[2pt]
 &=\frac{E^2(r-R)}{(r+R)^3}.
\end{align*}
Поскольку $E^2>0$ и $(r+R)^3>0$, знак производной определяется только знаком $r-R$.

\lessonstep{Шаг 3. Найти и доказать максимум}
\[
 P'(R)=0\iff r-R=0\iff R=r=0{,}8\ \text{Ом}.
\]
Для $0<R<r$ имеем $P'>0$, а для $R>r$ имеем $P'<0$. В частности, при $R=1$ Ом:
\[
 P'(1)=\frac{60^2(0{,}8-1)}{(0{,}8+1)^3}<0.
\]
\begin{center}
\begin{tikzpicture}[>=Latex,x=1cm,y=1cm,scale=1]
\draw[->,accent] (0,0)--(12,0) node[right] {$R$};
\draw[accent] (.9,-.1)--(.9,.1) node[below=5pt] {$0$};
\draw[accent] (5.6,-.1)--(5.6,.1) node[below=5pt] {$0{,}8$};
\node at (3,0.5) {$P'>0$};\node at (8.5,.5) {$P'<0$};
\draw[accent,thick,->] (2,-.65)--(5.1,-.2);
\draw[accent,thick,->] (6.1,-.2)--(10.4,-.65);
\end{tikzpicture}
\end{center}
\textbf{Вывод.} Мощность возрастает до $R=r$, затем убывает: это максимум.
\[\result{R=0{,}8\ \text{Ом}}\]
\clearpage

\lessonsection{3. Максимум распределения скоростей молекул}
\textbf{Условие.} Для кислорода распределение Максвелла описывается функцией
\[
 f(v)=4\pi v^2\left(\frac{M}{2\pi R_{\text{г}}T}\right)^{3/2}
       e^{-\frac{Mv^2}{2R_{\text{г}}T}},
\]
где $M=0{,}032$ кг/моль, $T=500$ К, $R_{\text{г}}=8{,}31$ Дж/(моль$\cdot$К). Найдите скорость $v$, при которой функция максимальна.

\lessonstep{Шаг 1. Скрыть постоянные в обозначениях}
\[
 A=4\pi\left(\frac{M}{2\pi R_{\text{г}}T}\right)^{3/2}>0,
 \qquad B=-\frac{M}{2R_{\text{г}}T}<0.
\]
Теперь длинная запись превращается в удобную функцию:
\[
 \result{f(v)=Av^2e^{Bv^2}},\qquad v\ge 0.
\]
Скорость $v$ --- переменная. $A$ и $B$ --- постоянные; при обратной подстановке важно сохранить минус в $B$.

\lessonstep{Шаг 2. Произведение и сложная функция}
\[
 (Av^2)'=2Av,\qquad (e^{Bv^2})'=2Bv e^{Bv^2}.
\]
\begin{align*}
 f'(v)&=2Av e^{Bv^2}+Av^2\cdot2Bv e^{Bv^2}\\
 &=\result{2Av e^{Bv^2}(1+Bv^2)}.
\end{align*}
При $v>0$ множители $2Av$ и $e^{Bv^2}$ положительны, поэтому знак $f'$ задаёт скобка $1+Bv^2$.

\lessonstep{Шаг 3. Критическая скорость}
\[
 1+Bv^2=0\iff v^2=-\frac1B=\frac{2R_{\text{г}}T}{M}.
\]
Выбираем положительный корень, поскольку скорость неотрицательна:
\[
 v_*=\sqrt{\frac{2R_{\text{г}}T}{M}}
    =\sqrt{\frac{2\cdot 8{,}31\cdot500}{0{,}032}}
    \approx\result{510\ \text{м/с}}.
\]
При $0<v<v_*$ выражение $1+Bv^2>0$, а при $v>v_*$ оно отрицательно. Следовательно, $f$ сначала возрастает, затем убывает. В точке $v_*$ достигается максимум плотности распределения.

\textbf{Замечание.} На границе $v=0$ имеем $f(0)=0$. Формула даёт наиболее вероятную скорость молекул; её не следует путать со средней или среднеквадратичной скоростью.
\clearpage

\lessonsection{4. Минимальный снос лодки течением}
\textbf{Условие.} Ширина реки $h=20$ м, скорость течения $u=3$ м/с, скорость лодки относительно воды $v=2{,}4$ м/с. Лодка направлена под углом $\alpha$ к направлению течения. Найдите наименьшее расстояние сноса $l$ до противоположного берега, если
\[
 l(\alpha)=\frac{h(u+v\cos\alpha)}{v\sin\alpha}.
\]

\lessonstep{Шаг 1. ОДЗ и выбор переменной}
Переменная --- угол $\alpha$. Чтобы лодка пересекала реку, необходимо $\sin\alpha>0$, то есть $0<\alpha<\pi$. Здесь $u>v>0$, поэтому $u+v\cos\alpha\ge u-v>0$; снос положителен.

\lessonstep{Шаг 2. Производная дроби}
\[
 U=h(u+v\cos\alpha),\quad U'=-hv\sin\alpha,
 \quad V=v\sin\alpha,\quad V'=v\cos\alpha.
\]
\begin{align*}
 l'(\alpha)&=\frac{-hv^2\sin^2\alpha-hv(u+v\cos\alpha)\cos\alpha}{v^2\sin^2\alpha}\\[2pt]
 &=\frac{-hv\bigl(v(\sin^2\alpha+\cos^2\alpha)+u\cos\alpha\bigr)}{v^2\sin^2\alpha}\\[2pt]
 &=\result{-\frac{h(v+u\cos\alpha)}{v\sin^2\alpha}}.
\end{align*}
Ключевое упрощение: $\sin^2\alpha+\cos^2\alpha=1$.

\lessonstep{Шаг 3. Критический угол и значение $l$}
Поскольку знаменатель положителен, приравниваем к нулю числитель:
\[
 v+u\cos\alpha=0\quad\Longrightarrow\quad
 \cos\alpha=-\frac vu=-\frac{2{,}4}{3}=-0{,}8.
\]
В интервале $(0,\pi)$ это тупой угол, поэтому
\[
 \sin\alpha=\sqrt{1-(-0{,}8)^2}=0{,}6.
\]
Тогда
\[
 l_{\min}=\frac{20(3+2{,}4\cdot(-0{,}8))}{2{,}4\cdot0{,}6}
 =\frac{20\cdot1{,}08}{1{,}44}=\result{15\ \text{м}}.
\]
\clearpage

\lessonsection{5. Как строго доказать минимум сноса}
Обозначим через $\alpha_*=\arccos(-0{,}8)$ найденный угол. При $0<\alpha<\alpha_*$ косинус больше $-0{,}8$, поэтому $v+u\cos\alpha>0$, откуда $l'(\alpha)<0$. При $\alpha_*<\alpha<\pi$ косинус меньше $-0{,}8$, поэтому $l'(\alpha)>0$.

\begin{center}
\begin{tikzpicture}[>=Latex,x=1cm,y=1cm]
\draw[->,accent] (0,0)--(12,0) node[right] {$\alpha$};
\draw[accent] (1,-.1)--(1,.1) node[below=5pt] {$0$};
\draw[accent] (6,-.1)--(6,.1) node[below=5pt] {$\alpha_*$};
\draw[accent] (11,-.1)--(11,.1) node[below=5pt] {$\pi$};
\node at (3.5,.5) {$l'<0$};\node at (8.5,.5) {$l'>0$};
\draw[accent,thick,->] (2,-.2)--(5.5,-.65);
\draw[accent,thick,->] (6.5,-.65)--(10,-.2);
\end{tikzpicture}
\end{center}

\textbf{Вывод.} До $\alpha_*$ расстояние сноса уменьшается, после $\alpha_*$ увеличивается. Следовательно, $l_{\min}=15$ м. На границах $\alpha\to0+$ и $\alpha\to\pi-$ расстояние растёт без ограничения, так как $u>v$.

\lessonstep{Почему важно правильно подписывать ось}
Можно ввести новую переменную $x=\cos\alpha\in(-1,1)$. Тогда
\[
 l(x)=\frac{h(u+vx)}{v\sqrt{1-x^2}},\qquad
 \frac{dl}{dx}=\frac{h(v+ux)}{v(1-x^2)^{3/2}}.
\]
Для $x<-0{,}8$ производная по $x$ отрицательна, для $x>-0{,}8$ --- положительна, что тоже доказывает минимум.

\textbf{Но:} $l'(\alpha)$ и $\dfrac{dl}{dx}$ --- производные по разным переменным. Их знаки нельзя механически переносить на одну и ту же ось: косинус убывает при возрастании угла $\alpha$ на $(0,\pi)$.

\lessonsection{6. Типичные ошибки}
\begin{itemize}
\item \textbf{Производная найдена, а экстремум не доказан.} Нужна таблица знаков или равноценное обоснование монотонности.
\item \textbf{Степень в знаменателе потеряна.} В электрической задаче после дифференцирования появляется $(r+R)^4$, затем знаменатель сокращается до $(r+R)^3$.
\item \textbf{Неверно обозначена константа.} В распределении Максвелла $B=-M/(2R_{\text{г}}T)$ обязательно отрицательна.
\item \textbf{Пропущено ограничение.} $R>0$, $v\ge0$, $0<\alpha<\pi$; физически недопустимые решения исключают.
\item \textbf{Перепутана переменная в таблице знаков.} Для оси $\alpha$ используйте $l'(\alpha)$; для оси $x=\cos\alpha$ --- производную $dl/dx$.
\end{itemize}
\clearpage

\lessonsection{7. Итоговый чек-лист}
Перед сдачей задачи на экстремум убедитесь, что каждый пункт выполнен.
\begin{itemize}[label=$\square$,itemsep=8pt]
\item Я определил, какую величину нужно максимизировать или минимизировать.
\item Я выбрал единственную переменную и отделил её от постоянных.
\item Я записал все математические и физические ограничения.
\item Я аккуратно вывел функцию и правильно вычислил её производную.
\item Я упростил производную до выражения с понятным знаком.
\item Я нашёл все допустимые критические точки и рассмотрел необходимые границы.
\item Я показал знаки производной и направление изменения функции.
\item Я обосновал максимум или минимум, подставил найденное значение в исходную формулу и указал единицы измерения.
\end{itemize}

\vspace{5mm}
\lessonstep{Сводка трёх моделей}
\renewcommand{\arraystretch}{1.45}
\begin{tabularx}{\linewidth}{@{}p{.26\linewidth}X p{.31\linewidth}@{}}
\toprule
Модель & Целевая функция & Результат\\
\midrule
Электрическая цепь & $P(R)=\dfrac{E^2R}{(r+R)^2}$ & $R=r=0{,}8$ Ом (максимум)\\[3pt]
Распределение Максвелла & $f(v)=Av^2e^{Bv^2}$ & $v=\sqrt{\dfrac{2R_{\text{г}}T}{M}}\approx510$ м/с (максимум)\\[3pt]
Снос лодки & $l(\alpha)=\dfrac{h(u+v\cos\alpha)}{v\sin\alpha}$ & $\cos\alpha=-0{,}8$, $l=15$ м (минимум)\\
\bottomrule
\end{tabularx}

\vspace{8mm}
\textbf{Главная идея урока.} Сложная физическая формула становится управляемой, если сначала выбрать переменную, сгруппировать константы, затем найти производную и строго проверить её знак. Именно это завершает доказательство экстремума.

\end{document}
